\section{Routing and fusion: protocols and supplementary results}
\label{app:assembly_results}
% Non-destructive TeX annotation layer. Scientific pixels, legends and
% numerical error annotations remain in the original archived PNGs.
\newcommand{\FieldPanel}[2]{%
\begin{tikzpicture}
\node[anchor=south west,inner sep=0] (im) at (0,0)
 {\includegraphics[width=0.80\linewidth]{#1}};
\begin{scope}[x={(im.south east)},y={(im.north west)}]
 \fill[white] (0,0.13) rectangle (0.055,0.98);
 \node[rotate=90,font=\scriptsize\bfseries] at (0.023,0.863) {Reference};
 \node[rotate=90,font=\scriptsize\bfseries] at (0.023,0.656) {#2};
 \node[rotate=90,font=\scriptsize\bfseries] at (0.023,0.449) {Hopf};
 \node[rotate=90,font=\scriptsize\bfseries] at (0.023,0.242) {RAL-MoE-ROM};
\end{scope}
\end{tikzpicture}%
}
\newcommand{\ErrorPanel}[2]{%
\begin{tikzpicture}
\node[anchor=south west,inner sep=0] (im) at (0,0)
 {\includegraphics[width=0.80\linewidth]{#1}};
\begin{scope}[x={(im.south east)},y={(im.north west)}]
 \fill[white] (0,0.17) rectangle (0.055,0.985);
 \node[rotate=90,font=\scriptsize\bfseries] at (0.023,0.83) {#2};
 \node[rotate=90,font=\scriptsize\bfseries] at (0.023,0.57) {Hopf};
 \node[rotate=90,font=\scriptsize\bfseries] at (0.023,0.31) {RAL-MoE-ROM};
 \foreach \x/\y/\lab in {0.080/0.932/i,0.583/0.932/j,0.080/0.673/k,0.583/0.673/l,0.080/0.415/m,0.583/0.415/n}{
  \fill[white] (\x-0.018,\y-0.020) rectangle (\x+0.018,\y+0.020);
  \node[font=\scriptsize\bfseries,inner sep=0] at (\x,\y) {(\lab)};
 }
\end{scope}
\end{tikzpicture}%
}


\subsection{Square: full-window fusion protocol}
\label{app:square}

The primary assembly study fixes the local predictors and E2 router.
The S--H test set uses $Re=95.1,95.3$, and the H--P set uses
$Re=100.5,102.0$, with eight windows per parameter and $K=24$.
Only the prescribed adjacent pair is compared at each interface.
The main text reports full-window joint errors; the following
subsections supply component errors, control fitting, and secondary
terminal-step measurements, rather than additional independent trials.

\paragraph{Scalar controls and component errors.}
Table~\ref{tab:t2c_full_window_controls} uses all $K=24$ steps, first
averaging windows within each Reynolds number and then weighting the
two held-out Reynolds numbers equally. It also reports the dimensionless
squared relative objective, $L_{\mathrm{sq}}=\operatorname{mean}
[(E_u/100)^2+(E_p/100)^2]$. This is different from the displayed
percentage joint error. The constant-$\alpha$ control minimizes the
training-window squared objective analytically over $[0,1]$; its weights
are $0.357850555$ in $S$--$H$ and $0.928683495$ in $H$--$P$.
The constant-logit control fits a single training-only offset in
$[-20,20]$, giving $0.490963974$ and $3.480802826$, respectively.
The validation-fixed candidate is selected on validation squared error,
rather than chosen a posteriori on held-out data. Both candidates are
finite and no held-out windows are excluded in this reevaluation.
These are deterministic scalar controls for the archived candidates.

\begin{table}[H]
\centering
\caption{Centered-square full-window $K=24$ reevaluation using the
archived physical-field quadratic cache. Errors are percentages against
the cached POD reference; the last column is $10^4 L_{\mathrm{sq}}$.
This table uses a single frozen gate checkpoint, not the three-seed mean.
Joint errors appear only in Table~\ref{tab:routing_fusion}.}
\label{tab:t2c_full_window_controls}
\small
\setlength{\tabcolsep}{3.5pt}
\begin{tabular}{@{}llrrr@{}}
\toprule
Overlap & Method & $E_u$ & $E_p$ & $10^4L_{\mathrm{sq}}$\\
\midrule
$S$--$H$ & E2 Top-1 & 0.5355 & 1.0432 & 5.6364 \\
 & E2 probability & 0.6624 & 1.1521 & 4.3056 \\
 & Equal blend & 1.0133 & 1.3717 & 4.3665 \\
 & Train-fixed $\alpha$ & 0.8128 & 1.2462 & 4.1324 \\
 & Train-fixed logit offset & 0.7857 & 1.2292 & 4.1398 \\
 & Validation-fixed candidate & 0.5355 & 1.0432 & 5.6364 \\
 & T2-C & 0.3831 & 0.7123 & 2.3199 \\
\midrule
$H$--$P$ & E2 Top-1 & 0.7263 & 0.9161 & 4.8921 \\
 & E2 probability & 0.6108 & 0.8381 & 2.9087 \\
 & Equal blend & 0.5920 & 0.8010 & 2.5103 \\
 & Train-fixed $\alpha$ & 0.5016 & 0.4752 & 0.5818 \\
 & Train-fixed logit offset & 0.5134 & 0.4687 & 0.5941 \\
 & Validation-fixed candidate & 0.5310 & 0.4758 & 0.6323 \\
 & T2-C & 0.4263 & 0.3860 & 0.4072 \\
\bottomrule
\end{tabular}
\end{table}

\paragraph{Gate-only repeated training.}
Seeds 42001, 42002, and 42003 retrain only the T2-C gate.
The specialists, E2, candidate trajectories, data splits, and
normalization remain fixed. Full T2-C and the architecture-matched
$\mu$-only control use the same optimization and selection protocol
(Appendix~\ref{app:implementation}); the latter zeros standardized
history channels. Full-window joint means and sample standard deviations
are reported in Section~\ref{sec:exp_fusion} and
Figure~\ref{fig:square_overlap}, and are not tabulated again here.
A training-fitted logit offset can change the weight through the
parameter-dependent E2 preference but cannot distinguish histories
at the same parameter. T2-C also uses one constant weight per window;
it is not updated from the predicted trajectory.

\subsection{Square: terminal-step and field supplements}
\label{app:square_terminal}

This subsection complements the full-window assembly comparison in
Section~\ref{sec:exp_fusion}. We first examine accuracy at the final
prediction step, then resolve the results by Reynolds number and
gate-training seed, and finally present representative reconstructed
fields. The numerical tables use the same Square test windows and
candidate predictors as the main-text study, rather than an additional
independent test set.

\paragraph{Terminal-step comparison.}
Table~\ref{tab:square_fusion_full} reports velocity, pressure, and joint
errors at $k=K=24$, averaged over the 16 held-out windows in each overlap.
These endpoint errors complement, but are distinct from, averages over
all 24 prediction steps. The gate entries in this table correspond
to individual validation-selected checkpoints; repeated gate training
is summarized separately in Table~\ref{tab:t2c_seed_variability}.

T2-C has the lowest terminal velocity, pressure, and joint errors among
the deployable methods listed in Table~\ref{tab:square_fusion_full}.
The more accurate individual candidate differs between overlaps:
Hopf is preferable in S--H, whereas Periodic is preferable in H--P.
T2-C improves on both of these single-candidate references, indicating
that its benefit is not limited to correcting the E2 Top-1 choice.
The improvement observed in the full-window study is therefore also
present at the evaluated endpoint; this does not establish accuracy
beyond the tested prediction horizon.

The convex oracle provides an a posteriori reference using the true
fields. Its weight minimizes a full-window objective and is evaluated
here at the terminal step. It is neither deployable nor a guaranteed
lower bound for each terminal error component.

\begin{table}[H]
\centering
\caption{
Centered-square terminal-step errors at $k=K=24$ (\%). The convex oracle is an a posteriori reference unavailable during prediction.
T2-C $\mu$-only uses the same gate architecture as T2-C but removes the standardized history descriptor. Bold indicates the best deployable method within each overlap.
}
\label{tab:square_fusion_full}
\small
\setlength{\tabcolsep}{5pt}
\begin{tabular}{@{}llrrr@{}}
\toprule
Overlap & Predictor & $E_u$ & $E_p$ & $E_{\mathrm{joint}}$ \\
\midrule
$S$--$H$ & Steady specialist & 2.5976 & 2.2450 & 4.8426 \\
 & Hopf specialist & 0.4942 & 1.2356 & 1.7298 \\
 & E2 Top-1 & 0.4942 & 1.2356 & 1.7298 \\
 & Equal blend & 1.3644 & 1.6677 & 3.0321 \\
 & E2-probability blend & 0.7898 & 1.3949 & 2.1847 \\
 & T2-C $\mu$-only & 0.7898 & 1.3948 & 2.1846 \\
 & T2-C & \textbf{0.4615} & \textbf{0.7343} & \textbf{1.1958} \\
 & Convex oracle & 0.4659 & 0.7183 & 1.1842 \\
\midrule
$H$--$P$ & Periodic specialist & 0.6420 & 0.5659 & 1.2079 \\
 & Hopf specialist & 1.0034 & 1.8684 & 2.8718 \\
 & E2 Top-1 & 0.7422 & 1.1052 & 1.8473 \\
 & Equal blend & 0.6077 & 1.0069 & 1.6146 \\
 & E2-probability blend & 0.6154 & 1.0464 & 1.6618 \\
 & T2-C $\mu$-only & 0.5875 & 0.5629 & 1.1504 \\
 & T2-C & \textbf{0.4509} & \textbf{0.3956} & \textbf{0.8465} \\
 & Convex oracle & 0.4479 & 0.3905 & 0.8384 \\
\bottomrule
\end{tabular}
\end{table}

\paragraph{Parameter-resolved behavior.}
Table~\ref{tab:square_parameter} separates the endpoint comparison
by Reynolds number. T2-C reduces joint error relative to E2 Top-1
at all four tested parameter values, so the aggregate improvement
is not confined to a single parameter.
In S--H, E2 selects Hopf at both Reynolds numbers, while T2-C assigns
approximately $0.097$ mean weight to Steady and retains a predominantly
Hopf prediction. In H--P, E2 switches from Hopf to Periodic as the
parameter changes; the mean Periodic fusion weight increases from
$0.5870$ to $0.7094$.
These averages describe how the assembled predictions differ across
the tested parameters. They do not, by themselves, isolate dependence
on history at a fixed parameter; that contribution is examined by
the matched history ablation.

\begin{table}[H]
\centering
\caption{
Parameter-wise centered-square routing and T2-C terminal-step results at $k=K=24$.
Errors are percentages; reduction is
$100(E_{\mathrm{joint}}^{\mathrm{E2}}-
E_{\mathrm{joint}}^{\mathrm{T2-C}})/E_{\mathrm{joint}}^{\mathrm{E2}}$.
Each row averages eight held-out windows.
The mean weight is the Steady-specialist weight for $S$--$H$
and the Periodic-specialist weight for $H$--$P$.
}
\label{tab:square_parameter}
\small
\setlength{\tabcolsep}{3.5pt}
\begin{tabular}{lrlrrrr}
\toprule
Overlap & $Re$ & E2 choice & $\bar\alpha$ & E2 $E_{\mathrm{joint}}$ & T2-C $E_{\mathrm{joint}}$ & Reduction \\
\midrule
$S$--$H$ & 95.1 & Hopf & 0.0967 & 1.7140 & 1.1827 & 31.00\% \\
$S$--$H$ & 95.3 & Hopf & 0.0971 & 1.7456 & 1.2090 & 30.74\% \\
$H$--$P$ & 100.5 & Hopf & 0.5870 & 2.6485 & 0.8571 & 67.64\% \\
$H$--$P$ & 102.0 & Periodic & 0.7094 & 1.0462 & 0.8360 & 20.09\% \\
\bottomrule
\end{tabular}
\end{table}

\paragraph{Terminal-step gate variability.}
Table~\ref{tab:t2c_seed_variability} evaluates the endpoint errors
of the same three paired gate fits used in the main-text history
ablation (seeds 42001, 42002, and 42003).
Only the gate is retrained; local predictors, E2, candidate trajectories,
data splits, and normalization are fixed.
The $\mu$-only control retains the parameter and E2-preference inputs
but zeros the standardized history channels without changing the
gate architecture.

Full T2-C has lower mean velocity and pressure errors than the
history-ablated control in both overlaps. Thus the benefit of history
conditioning is visible at the endpoint as well as in the full-window
metric. The sample standard deviations measure sensitivity to
gate-training seeds, not variation across Reynolds numbers, test-set
uncertainty, or the variability of retraining the full hierarchy.

\begin{table}[H]
\centering
\caption{
Three-seed variability of the centered-square T2-C gate at the $K=24$
terminal step. Errors (\%) are mean $\pm$ sample standard deviation
over three gate-training seeds, after averaging the held-out windows
for each seed. Bold identifies the smaller mean within each overlap
and metric, not a statistical significance test.
}
\label{tab:t2c_seed_variability}
\small
\setlength{\tabcolsep}{4pt}
\begin{tabular}{@{}llccc@{}}
\toprule
Overlap & Gate & $E_u$ & $E_p$ & $E_{\mathrm{joint}}$ \\
\midrule
$S$--$H$ & T2-C $\mu$-only & $0.79006\pm0.00026$ & $1.39498\pm0.00012$ & $2.18504\pm0.00038$ \\
 & Full T2-C & $\mathbf{0.46156\pm0.00004}$ & $\mathbf{0.73416\pm0.00014}$ & $\mathbf{1.19572\pm0.00010}$ \\
\midrule
$H$--$P$ & T2-C $\mu$-only & $0.59067\pm0.00461$ & $0.56050\pm0.00435$ & $1.15116\pm0.00070$ \\
 & Full T2-C & $\mathbf{0.45211\pm0.00138}$ & $\mathbf{0.39793\pm0.00215}$ & $\mathbf{0.85004\pm0.00347}$ \\
\bottomrule
\end{tabular}
\end{table}

\paragraph{Representative field reconstructions.}
Figures~\ref{fig:app_square_sh_full} and \ref{fig:app_square_hp_full}
supplement the scalar errors with representative S--H and H--P fields
at $Re=95.1$ and $Re=100.5$, respectively.
Each figure juxtaposes the common reference, both local candidates,
and T2-C for velocity magnitude and gauge-corrected pressure, followed
by pointwise absolute-error maps. These views allow spatial discrepancies
to be inspected rather than reduced to a single aggregate norm.
The archived snapshots are qualitative supplements: they are not
identified here as the $k=24$ frames, and their visual agreement
does not replace the endpoint or full-window statistics.

\begin{figure}[p]
\centering
\FieldPanel{figures/original/SH_2.png}{Steady}

\vspace{0.35em}
\ErrorPanel{figures/original/SH_2_1.png}{Steady}
\caption{Representative held-out prediction in the centered-square
$S$--$H$ overlap at $Re=95.1$. The upper panel compares the archived reference reconstruction,
Steady specialist, Hopf specialist, and T2-C prediction for velocity
magnitude and gauge-corrected pressure. The lower panel gives the
corresponding pointwise absolute errors. Panel identifiers run (a)--(h) for fields and (i)--(n) for errors. The displayed snapshots and color scales are retained from the archived figures; row labels and identifiers are typeset separately.}
\label{fig:app_square_sh_full}
\end{figure}

\begin{figure}[p]
\centering
\FieldPanel{figures/original/HP_2.png}{Periodic}

\vspace{0.35em}
\ErrorPanel{figures/original/HP_2_1.png}{Periodic}
\caption{Representative held-out prediction in the centered-square
$H$--$P$ overlap at $Re=100.5$. The upper panel compares the archived reference reconstruction, Periodic specialist, Hopf specialist, and T2-C prediction for velocity magnitude and gauge-corrected pressure. The lower panel gives the
corresponding pointwise absolute errors. Panel identifiers run (a)--(h) for fields and (i)--(n) for errors. The displayed snapshots and color scales are retained from the archived figures; row labels and identifiers are typeset separately.}
\label{fig:app_square_hp_full}
\end{figure}

\FloatBarrier

\subsection{Pinball: additional assembly results}
\label{app:pinball_overlap_results}
This archived study compares S--H and H--P assembly at $K=24,56$.
The formal outer candidates are Deep-FNN-H3 local predictors, not
sparse-MoE specialists. Table~\ref{tab:pinball_overlap} retains the
archived joint-error summaries; it is not pooled with the Square
full-window ablation or interpreted as a replication of its scalar
controls and three gate seeds. Candidate identities and the oracle
convention are specified in the caption.

\begin{table}[H]
\centering
\caption{
Fluidic-pinball overlap joint errors (\%). For $S$--$H$, Candidate~1 is the Steady specialist and Candidate~2 is Hopf; for $H$--$P$, Candidate~1 is the Periodic specialist and Candidate~2 is Hopf.
The convex oracle selects its weight a posteriori for each evaluation window and is unavailable during prediction.
}
\label{tab:pinball_overlap}
\small
\setlength{\tabcolsep}{3.5pt}
\begin{tabular}{lrccccc}
\toprule
Overlap & $K$ & Cand.~1 & Cand.~2 & E2-probability blend & T2-C & Convex oracle \\
\midrule
$S$--$H$ & 24 & 0.6014 & 3.1925 & 1.7511 & 0.5897 & 0.5889 \\
$S$--$H$ & 56 & 0.6927 & 3.1193 & 1.7714 & 0.6767 & 0.6763 \\
$H$--$P$ & 24 & 0.0990 & 0.5146 & 0.2972 & 0.0945 & 0.0928 \\
$H$--$P$ & 56 & 0.1028 & 0.5346 & 0.3085 & 0.0978 & 0.0955 \\
\bottomrule
\end{tabular}
\end{table}
